\documentclass[border=4pt]{standalone}

\usepackage{tikz}

\definecolor{plnavy}{HTML}{1F4E89}
\definecolor{plrose}{HTML}{B85A52}
\definecolor{plsage}{HTML}{93CC9E}
\definecolor{plsand}{HTML}{E0BE78}
\definecolor{plplum}{HTML}{8E7BB5}
\definecolor{plgray}{HTML}{6E6E6E}

% ==== data: one entry per spoke, clockwise from the top ====
\def\radarn{10}
\def\radarnames{{Model A\\Set 1},{Model B\\Set 1},{Model C\\Set 1},{Model D\\Set 1},{Model A\\Set 2},{Model B\\Set 2},{Model C\\Set 2},{Model D\\Set 2},{Model A\\Set 3},{Model B\\Set 3}}
\def\axmin{{80,80,80,80,40,40,40,55,55,55}}
\def\axmax{{91,91,91,91,61,61,61,85,85,85}}
\def\seriesA{{83.2,84.6,85.0,85.2,50.3,51.0,58.4,72.5,70.1,79.5}}
\def\seriesB{{85.4,86.0,87.9,89.1,51.8,52.0,59.3,73.6,71.0,80.6}}
\def\seriesC{{87.9,88.5,90.0,90.2,49.8,52.8,60.2,74.2,71.5,81.1}}

\pgfmathtruncatemacro\lastaxis{\radarn-1}
\def\R{3.2}
\newcommand{\radarpt}[2]{({90-#2*360/\radarn}:{\R*(#1[#2]-\axmin[#2])/(\axmax[#2]-\axmin[#2])})}
\newcommand{\radarseries}[2]{
  \filldraw[draw=#2, fill=#2, fill opacity=0.07, line width=1pt, line join=round] \radarpt{#1}{0} \foreach \i in {1,...,\lastaxis}{ -- \radarpt{#1}{\i} } -- cycle;
  \foreach \i in {0,...,\lastaxis}{ \fill[#2] \radarpt{#1}{\i} circle[radius=1.2pt]; }
}

\begin{document}
\begin{tikzpicture}
  \foreach \level in {0.3333,0.6667,1}{
    \draw[black!70, line width=0.4pt, line join=round] (90:\level*\R) \foreach \i in {1,...,\lastaxis}{ -- ({90-\i*360/\radarn}:\level*\R) } -- cycle;
  }
  \radarseries{\seriesA}{plrose}
  \radarseries{\seriesB}{plsage!70!black}
  \radarseries{\seriesC}{plnavy}
  \foreach \name [count=\c from 0] in \radarnames {
    \pgfmathsetmacro\a{90-\c*360/\radarn}
    \draw[black!35, line width=0.3pt] (0,0) -- (\a:\R);
    \node[align=center, font=\scriptsize] at (\a:\R+0.85) {\name};
    \foreach \level in {0.3333,0.6667,1}{
      \pgfmathsetmacro\v{\axmin[\c]+\level*(\axmax[\c]-\axmin[\c])}
      \node[font=\tiny, fill=white, inner sep=0.5pt] at (\a:\level*\R) {\pgfmathprintnumber[fixed, precision=0]{\v}};
    }
  }
  \foreach \c/\l [count=\k] in {plrose/Baseline, plsage!70!black/Variant, plnavy/Ours}{
    \draw[\c, line width=1pt] (4.5,-2.5-0.42*\k) -- ++(0.6,0) node[right, black, font=\small] {\l};
  }
\end{tikzpicture}
\end{document}
